% Created 2026-07-17 Fri 18:13
% Intended LaTeX compiler: pdflatex
\documentclass[11pt]{article}

\usepackage[utf8]{inputenc}
\usepackage[T1]{fontenc}
\usepackage{graphicx}
\usepackage{longtable}
\usepackage{wrapfig}
\usepackage{rotating}
\usepackage[normalem]{ulem}
\usepackage{amsmath}
\usepackage{amssymb}
\usepackage{capt-of}
\usepackage{hyperref}
\usepackage{fullpage}
\author{Anup Zope}
\date{\today}
\title{C++ Text Template Engine}
\hypersetup{
 pdfauthor={Anup Zope},
 pdftitle={C++ Text Template Engine},
 pdfkeywords={},
 pdfsubject={},
 pdfcreator={},
 pdflang={English}}
\begin{document}

\maketitle
\tableofcontents

\section{Section}
\label{sec:org5374fa6}

A C++ text template engine with \texttt{\$\$variable\$\$} substitution. It is created to
streamline the process of generating biolerplace code required to generate CUDA
enabled Loci rules from Loci preprocessor.
\section{Features}
\label{sec:org29f0111}

\begin{itemize}
\item \texttt{\$\$variable\$\$} substitution from a hierarchical context
\item \texttt{\$\$user.name\$\$} nested key paths
\item \texttt{\$\$::path\$\$} root-scoped paths (always resolve from the render context root)
\item \texttt{\$\$\#if condition\$\$...\$\$/if\$\$} conditionals with \texttt{\$\$\#else if condition\$\$} and
optional \texttt{\$\$\#else\$\$}
\item Logical conditions: \texttt{\&\&}, \texttt{||}, \texttt{!}, and parentheses (\texttt{\&\&} binds tighter than
\texttt{||})
\item \texttt{\$\$\#each list\$\$...\$\$/each\$\$} loops over arrays (exposes \texttt{\$\$@index\$\$} and
\texttt{\$\$@size\$\$} per item)
\item Nested \texttt{\$\$\#if\$\$} blocks and \texttt{\$\$\#if\$\$} inside \texttt{\$\$\#each\$\$}
\item Named nested templates via \texttt{define()} and \texttt{\$\$> name\$\$} includes
\item \texttt{\#} and \texttt{/} prefixes reserved for control statements; unknown directives fail
at parse time with line and column
\item User-defined C++ callbacks for newline character in template text
\item Whitespace trimming inside placeholders
\item Optional behavior for missing variables (empty string or keep placeholder)
\item File-based template rendering
\end{itemize}
\section{Usage}
\label{sec:orge511e08}

\begin{verbatim}
#include "template.h"

Loci::TemplateEngine engine;

Loci::TemplateContext context;
context.set("greeting", "Hello");
context.set("name", "World");
context.set("active", "true");
context.set("count", 3);

std::string output = engine.render("$$greeting$$, $$name$$!", context);
// output: "Hello, World!"
\end{verbatim}

\texttt{set()} has typed overloads for \texttt{std::string}, \texttt{const char*}, \texttt{bool}, \texttt{char},
integer types, and floating-point types. Values are stored in the context
variant and formatted when rendered (for example, \texttt{bool} renders as \texttt{1} or \texttt{0},
floats use \texttt{std::to\_string}).

Build a flat string map quickly with \texttt{from\_flat}:

\begin{verbatim}
auto context = Loci::TemplateContext::from_flat({
  {"name", "World"},
  {"active", "true"},
});
\end{verbatim}
\subsection{Nested Keys}
\label{sec:orgdd570af}

\begin{verbatim}
Loci::TemplateContext user;
user.set("name", "Alice");

Loci::TemplateContext context;
context.set_object("user", user);

engine.render("Hello, $$user.name$$!", context);
// output: "Hello, Alice!"
\end{verbatim}

\texttt{set\_object()} and \texttt{set\_array()} store nested data by value (deep copy). Use
\texttt{set\_object} for objects and \texttt{set\_array} for lists of \texttt{TemplateContext} items.
\subsection{Root scope (\texttt{::})}
\label{sec:orgd17ee1c}

Prefix a path with \texttt{::} to resolve from the root context instead of the current
loop scope. Use \texttt{.} to descend into nested objects:

\begin{verbatim}
// Inside a loop where the item also has "label":
engine.render("$$#each rows$$$$label$$ / $$::label$$$$/each$$", context);
// item label, then root label
\end{verbatim}

Works in variables (\texttt{\$\$::user.name\$\$}), conditionals (\texttt{\$\$\#if ::active\$\$}), and
loops (\texttt{\$\$\#each ::items\$\$}).

Inside a loop, unqualified paths resolve from the current item first, then fall
back to root. Use \texttt{::} when a loop item shadows a root key.
\subsection{Conditionals}
\label{sec:org9ef13d3}

\begin{verbatim}
context.set("active", "true");
engine.render("$$#if active$$On$$#else$$Off$$/if$$", context);
// output: "On"

context.set("active", "false");
context.set("backup", "true");
engine.render("$$#if active$$On$$#else if backup$$Fallback$$#else$$Off$$/if$$", context);
// output: "Fallback"
\end{verbatim}

\texttt{\$\$\#elseif condition\$\$} is also accepted as an alias for \texttt{\$\$\#else if
condition\$\$}.

Conditions support logical operators with parentheses for grouping:

\begin{verbatim}
engine.render("$$#if active && backup$$On$$/if$$", context);
engine.render("$$#if ::published || draft$$Visible$$/if$$", context);
engine.render("$$#if !(archived || hidden)$$Show$$/if$$", context);
\end{verbatim}

Nested blocks are supported:

\begin{verbatim}
$$#if outer$$
  $$#if ::inner$$
    both true
  $$/if$$
$$/if$$
\end{verbatim}

Placeholders starting with \texttt{\#} or \texttt{/} are control statements only (\texttt{\#if},
\texttt{\#else}, \texttt{\#each}, \texttt{/if}, \texttt{/each}, and so on). Anything else is a parse error,
for example \texttt{\$\$/fi\$\$} reports \texttt{Unrecognized control statement '/fi' (did you
mean '/if'?)} with a line and column.

\texttt{is\_true} treats non-empty strings as true except \texttt{false}, \texttt{0}, \texttt{no}, \texttt{off}.
Non-empty arrays and objects are also true. Typed booleans and non-zero numbers
are true in conditions.
\subsection{Loops}
\label{sec:org9a9f6d3}

\begin{verbatim}
Loci::TemplateContext a;
a.set("name", "Alice");
Loci::TemplateContext b;
b.set("name", "Bob");

Loci::TemplateContext context;
context.set_array("users", {a, b});

engine.render("$$#each users$$$$@index$$: $$name$$, $$/each$$", context);
// output: "0: Alice, 1: Bob, "
\end{verbatim}

Inside a loop, fields from the current item are in scope. \texttt{\$\$@index\$\$} is the
zero-based index of the current item. \texttt{\$\$@size\$\$} is the number of items in the
array being iterated.

\texttt{\$\$\#if\$\$} may appear inside a loop body; conditions follow the same scope rules
as variables.
\subsection{Nested templates}
\label{sec:org18f0ee1}

Break large templates into named pieces with \texttt{define()}, then include them with
\texttt{\$\$> name\$\$}. Nested templates inherit the current render scope (including
\texttt{\#each} item scope).

\begin{verbatim}
Loci::TemplateEngine engine;

engine.define("status_line",
  "$$#if published$$Status: Published$$#else$$Status: Draft$$/if$$\n");

engine.define("task_row",
  "$$@index$$. $$name$$\n");

engine.define("task_list",
  "Tasks:\n"
  "$$#each tasks$$$$> task_row$$$$/each$$");

engine.define("report",
  "Report: $$title$$\n"
  "$$> status_line$$"
  "\n"
  "$$> task_list$$");

std::string output = engine.render_template("report", context);
\end{verbatim}

Includes are resolved at render time. Circular includes and unknown names throw.
Use \texttt{has\_template("name")} to check registration.
\subsubsection{Newline Callbacks}
\label{sec:org6e35c5e}

A named template can replace each literal newline as it is rendered:

\begin{verbatim}
engine.define(
  "declarations",
  "$$#each inputs$$$$ctype$$ $$vname$$;\n$$/each$$",
  [](const Loci::NewlineEvent& event) {
    // The current #each item is available through event.context.
    return "\n#line 42 \"rules.template\"\n";
  });
\end{verbatim}

The callback runs once for every literal \texttt{\textbackslash{}n} in text emitted by that
definition, including repeated loop bodies. Its return value replaces the
newline; return \texttt{"\textbackslash{}n"} to preserve it unchanged. Newlines originating in
variable values do not invoke the callback.

Each nested template uses its own callback. A nested template without one
emits ordinary newlines, and the caller's callback is restored when the
include returns.
\subsection{Render From File}
\label{sec:org4bdac99}

\begin{verbatim}
std::string output = engine.render_file("templates/report.txt", context);
\end{verbatim}
\subsection{Options}
\label{sec:orgbddae7e}

\begin{verbatim}
Loci::RenderOptions options;
options.keep_missing_placeholders = true;  // leave $$missing$$ unchanged

Loci::TemplateEngine engine(options);
\end{verbatim}
\end{document}
